\newcommand{\thesisTitleKorean}{Cross-Modal 지식 증류와 Hybrid CNN-Transformer를 활용한 향상된 의료 이미지 분할 기법}

{\centering \Large \bf \thesisTitleKorean \par} \vspace{1.0cm}

hhuhuhu T\textsubscript{1ce} MRI 시퀀스를 효과적으로 분할할 수 있도록 설계되었으며, 


T\textsubscript{1} MRI 스캔과 구조적 유사성을 공유하여 충분한 정보를 제공한다. 
프레임워크는 성과 중심 응답 기반 KD와 협력적 심층 감독 융합 학습(CDSFL)이라는 두 가지 고유한 지식 증류(KD) 전략을 통합한다. 성과 중심 응답 기반 KD 전략은 각 교사 모델의 성과에 따라 신뢰도 가중치를 조정하여 KD 프로세스를 성과에 맞게 조정한다. 또한, CDSFL 모듈은 상호 학습을 촉진하여 다중 교사 모델의 학습 능력을 강화한다. 이러한 모델들의 결합된 지식은 학생 모델에 증류되어 딥러닝 감독을 향상시킨다. BraTS 데이터 세트를 통한 포괄적인 테스트 결과, 우리 프레임워크는 T\textsubscript{1ce} 및 T\textsubscript{1} 양식의 단일 모드 분할에서 기존의 최신 기술을 능가하는 유망한 결과를 보였다.